% ============================================================
% Chapter 5 — Electrical Design & Complete Wiring
% Source: PSYCON Engineering Design Document, Vol. 1, Ch. 5
% ============================================================
\section{Electrical Design \& Complete Wiring}
\label{sec:electrical-design}

\subsection{Electrical Design Philosophy}
\label{sec:electrical-design-philosophy}

The electrical system is designed around two independent ESP32-based modules: the \textbf{Wrist Module} (physiological sensing) and the \textbf{Audio Module} (speech sensing). Each module has an independent battery, independent charger, independent ESP32, and independent firmware. This separation minimizes cross-module interference, improves fault tolerance, and simplifies debugging.

\subsection{Power Architecture}
\label{sec:power-architecture-ch5}

\Cref{fig:power-architecture-detailed} shows the complete power chain for both modules, from battery cells through regulation to the connected sensor loads.

\begin{figure*}[t]
\centering
\begin{tikzpicture}[
    font=\small,
    pwrstage/.style={draw, rounded corners=3pt, minimum width=5.4cm, minimum height=0.85cm, align=center, line width=0.6pt, fill=green!8},
    ctrlstage/.style={draw, rounded corners=3pt, minimum width=5.4cm, minimum height=0.85cm, align=center, line width=0.6pt, fill=orange!10},
    loadstage/.style={draw, rounded corners=3pt, minimum width=5.4cm, minimum height=1.0cm, align=center, line width=0.6pt, fill=blue!8},
    arrow/.style={-{Stealth[length=2.2mm]}, line width=0.6pt}
]

% Wrist Module power chain
\node[pwrstage] (w1) at (0,5.2) {2$\times$500\,mAh LiPo (parallel)};
\node[pwrstage, below=5mm of w1] (w2) {TP4056 Charger + Protection};
\node[ctrlstage, below=5mm of w2] (w3) {ESP32 VIN / 5\,V Input\textsuperscript{*}};
\node[ctrlstage, below=5mm of w3] (w4) {3.3\,V Regulator};
\node[loadstage, below=5mm of w4] (w5) {MAX30102 \textbar\ MPU6050 \textbar\ MCP9808 \textbar\ ADS1115 \textbar\ GSR Front-End};

\draw[arrow] (w1.south) -- (w2.north);
\draw[arrow] (w2.south) -- (w3.north);
\draw[arrow] (w3.south) -- (w4.north);
\draw[arrow] (w4.south) -- (w5.north);

\node[above=2mm of w1, font=\bfseries\small] {Wrist Module};

% Audio Module power chain
\node[pwrstage] (a1) at (7.8,5.2) {500\,mAh LiPo};
\node[pwrstage, below=5mm of a1] (a2) {TP4056 Charger + Protection};
\node[ctrlstage, below=5mm of a2] (a3) {ESP32};
\node[loadstage, below=17.7mm of a3] (a4) {INMP441 \textbar\ TEMT6000};

\draw[arrow] (a1.south) -- (a2.north);
\draw[arrow] (a2.south) -- (a3.north);
\draw[arrow] (a3.south) -- (a4.north);

\node[above=2mm of a1, font=\bfseries\small] {Audio Module};

\end{tikzpicture}
\caption{Complete power distribution for the Wrist Module and Audio Module, from LiPo cells through the TP4056 charger/protection stage, ESP32 power input, regulation, and final sensor loads. \textsuperscript{*}Final power wiring depends on the exact ESP32 board; see the note below.}
\label{fig:power-architecture-detailed}
\end{figure*}

\begin{warningbox}
Many ESP32 DevKit boards expect a regulated 5\,V on \texttt{VIN} or USB, not a raw LiPo voltage. If powering directly from a LiPo cell, an appropriate regulator or a board designed for single-cell LiPo input is required. This must be verified on the specific board used.
\end{warningbox}

\subsection{Wrist Module GPIO Assignment}
\label{sec:wrist-gpio-assignment}

\Cref{tab:wrist-gpio} lists the GPIO assignment for the Wrist Module.

\begin{table}[H]
\centering
\caption{Wrist Module ESP32 GPIO assignment.}
\label{tab:wrist-gpio}
\small
\begin{tabularx}{\columnwidth}{@{}l Y@{}}
\toprule
\textbf{ESP32 GPIO} & \textbf{Device / Function} \\
\midrule
GPIO21 & I\textsuperscript{2}C SDA --- MAX30102, MPU6050, ADS1115, MCP9808 \\
GPIO22 & I\textsuperscript{2}C SCL --- MAX30102, MPU6050, ADS1115, MCP9808 \\
GPIO34 & ADS1115 \texttt{ALERT} (optional interrupt) \\
GPIO35 & Reserved \\
GPIO32 & Future battery monitor \\
GPIO33 & Future interrupt \\
\bottomrule
\end{tabularx}
\end{table}

\textbf{Power.} All Wrist Module sensors (MAX30102, ADS1115, MCP9808, MPU6050) are supplied from the 3.3\,V rail.

\textbf{Ground.} A common ground is shared by every module.

\subsection{Audio Module GPIO Assignment}
\label{sec:audio-gpio-assignment}

\Cref{tab:audio-gpio} lists the GPIO assignment for the Audio Module.

\begin{table}[H]
\centering
\caption{Audio Module ESP32 GPIO assignment.}
\label{tab:audio-gpio}
\small
\begin{tabularx}{\columnwidth}{@{}l Y@{}}
\toprule
\textbf{ESP32 GPIO} & \textbf{Device / Function} \\
\midrule
GPIO26 & I\textsuperscript{2}S BCLK \\
GPIO25 & I\textsuperscript{2}S LRCLK (WS) \\
GPIO33 & I\textsuperscript{2}S DATA \\
GPIO32 & TEMT6000 analog output \\
3.3\,V & Sensor VCC \\
GND & Common ground \\
\bottomrule
\end{tabularx}
\end{table}

\subsection{I\texorpdfstring{\textsuperscript{2}}{2}C Bus}
\label{sec:i2c-bus-ch5}

\Cref{fig:i2c-topology} shows the shared I\textsuperscript{2}C bus topology on the Wrist Module. \Cref{tab:i2c-addresses-ch5} lists the expected device addresses; no address conflicts are expected.

\begin{figure}[H]
\centering
\begin{tikzpicture}[
    font=\small,
    ctrl/.style={draw, rounded corners=3pt, minimum width=2.4cm, minimum height=0.8cm, align=center, line width=0.6pt, fill=orange!10},
    dev/.style={draw, rounded corners=3pt, minimum width=2.4cm, minimum height=0.7cm, align=center, line width=0.6pt, fill=gray!8},
    arrow/.style={-{Stealth[length=2mm]}, line width=0.5pt}
]
\node[ctrl] (esp) at (0,0) {ESP32\\ GPIO21 (SDA) / GPIO22 (SCL)};

\node[dev] (max) at (3.4,1.6) {MAX30102};
\node[dev] (ads) at (3.4,0.55) {ADS1115};
\node[dev] (mcp) at (3.4,-0.55) {MCP9808};
\node[dev] (mpu) at (3.4,-1.6) {MPU6050};

\draw[arrow] (esp.east) -- (max.west);
\draw[arrow] (esp.east) -- (ads.west);
\draw[arrow] (esp.east) -- (mcp.west);
\draw[arrow] (esp.east) -- (mpu.west);
\end{tikzpicture}
\caption{Shared I\textsuperscript{2}C bus topology on the Wrist Module: the ESP32 (GPIO21/SDA, GPIO22/SCL) drives MAX30102, ADS1115, MCP9808, and MPU6050 on a single bus.}
\label{fig:i2c-topology}
\end{figure}

\begin{table}[H]
\centering
\caption{Expected I\textsuperscript{2}C device addresses on the Wrist Module bus.}
\label{tab:i2c-addresses-ch5}
\small
\begin{tabularx}{\columnwidth}{@{}Y l@{}}
\toprule
\textbf{Device} & \textbf{Address} \\
\midrule
MCP9808 & \texttt{0x18} \\
ADS1115 & \texttt{0x48} \\
MAX30102 & \texttt{0x57} \\
MPU6050 & \texttt{0x68} \\
\bottomrule
\end{tabularx}
\end{table}

\subsection{I\texorpdfstring{\textsuperscript{2}}{2}S Bus}
\label{sec:i2s-bus-ch5}

The Audio Module I\textsuperscript{2}S bus connects the ESP32 to the INMP441 digital microphone as follows: GPIO26 $\rightarrow$ BCLK, GPIO25 $\rightarrow$ WS, GPIO33 $\rightarrow$ SD, 3.3\,V $\rightarrow$ VDD, GND $\rightarrow$ GND. The INMP441 L/R select pin is connected to GND, selecting the default left channel.

\subsection{ADS1115 Connections}
\label{sec:ads1115-connections}

\Cref{tab:ads1115-wiring} lists the complete ADS1115 wiring.

\begin{table}[H]
\centering
\caption{ADS1115 pin connections.}
\label{tab:ads1115-wiring}
\small
\begin{tabularx}{\columnwidth}{@{}l Y@{}}
\toprule
\textbf{ADS1115 Pin} & \textbf{Connection} \\
\midrule
VDD & 3.3\,V \\
GND & GND \\
SDA & GPIO21 \\
SCL & GPIO22 \\
ADDR & GND (address \texttt{0x48}) \\
ALRT & Optional, GPIO34 \\
A0 & GSR front end \\
A1--A3 & Reserved \\
\bottomrule
\end{tabularx}
\end{table}

\subsection{MAX30102 Connections}
\label{sec:max30102-connections}

\Cref{tab:max30102-wiring} lists the complete MAX30102 wiring.

\begin{table}[H]
\centering
\caption{MAX30102 pin connections.}
\label{tab:max30102-wiring}
\small
\begin{tabularx}{\columnwidth}{@{}l Y@{}}
\toprule
\textbf{Sensor Pin} & \textbf{ESP32 Connection} \\
\midrule
VIN & 3.3\,V \\
GND & GND \\
SDA & GPIO21 \\
SCL & GPIO22 \\
INT & Optional, GPIO27 \\
\bottomrule
\end{tabularx}
\end{table}

\subsection{MPU6050 Connections}
\label{sec:mpu6050-connections}

\Cref{tab:mpu6050-wiring} lists the complete MPU6050 wiring. The device address is \texttt{0x68}.

\begin{table}[H]
\centering
\caption{MPU6050 pin connections.}
\label{tab:mpu6050-wiring}
\small
\begin{tabularx}{\columnwidth}{@{}l Y@{}}
\toprule
\textbf{Pin} & \textbf{ESP32 Connection} \\
\midrule
VCC & 3.3\,V \\
GND & GND \\
SDA & GPIO21 \\
SCL & GPIO22 \\
INT & Optional, GPIO14 \\
AD0 & GND \\
\bottomrule
\end{tabularx}
\end{table}

\subsection{MCP9808 Connections}
\label{sec:mcp9808-connections}

\Cref{tab:mcp9808-wiring} lists the complete MCP9808 wiring. The device address is \texttt{0x18}.

\begin{table}[H]
\centering
\caption{MCP9808 pin connections.}
\label{tab:mcp9808-wiring}
\small
\begin{tabularx}{\columnwidth}{@{}l Y@{}}
\toprule
\textbf{Pin} & \textbf{ESP32 Connection} \\
\midrule
VIN & 3.3\,V \\
GND & GND \\
SDA & GPIO21 \\
SCL & GPIO22 \\
ALERT & Optional, GPIO13 \\
A0--A2 & GND \\
\bottomrule
\end{tabularx}
\end{table}

\subsection{TEMT6000 Connections}
\label{sec:temt6000-connections}

\Cref{tab:temt6000-wiring} lists the complete TEMT6000 wiring.

\begin{table}[H]
\centering
\caption{TEMT6000 pin connections.}
\label{tab:temt6000-wiring}
\small
\begin{tabularx}{\columnwidth}{@{}l Y@{}}
\toprule
\textbf{Pin} & \textbf{ESP32 Connection} \\
\midrule
VCC & 3.3\,V \\
GND & GND \\
SIG & GPIO32 (ADC) \\
\bottomrule
\end{tabularx}
\end{table}

\subsection{GSR Front End (Current Prototype)}
\label{sec:gsr-frontend-ch5}

\Cref{fig:gsr-frontend-block} shows the current prototype signal chain for electrodermal activity acquisition.

\begin{figure}[H]
\centering
\begin{tikzpicture}[
    font=\small,
    node distance=0mm,
    field/.style={draw, rounded corners=2pt, fill=blue!6, minimum width=3.2cm, minimum height=0.65cm, align=center, line width=0.5pt}
]
\node[field] (elec) {Electrode};
\node[field, below=6mm of elec] (cond) {Signal Conditioning};
\node[field, below=6mm of cond] (ads) {ADS1115 (A0)};
\node[field, below=6mm of ads] (esp) {ESP32};

\draw[-{Stealth[length=2mm]}, line width=0.6pt] (elec.south) -- (cond.north);
\draw[-{Stealth[length=2mm]}, line width=0.6pt] (cond.south) -- (ads.north);
\draw[-{Stealth[length=2mm]}, line width=0.6pt] (ads.south) -- (esp.north);
\end{tikzpicture}
\caption{Current prototype GSR signal chain: electrode $\rightarrow$ signal conditioning $\rightarrow$ ADS1115 channel A0 $\rightarrow$ ESP32.}
\label{fig:gsr-frontend-block}
\end{figure}

\begin{requirementbox}
\textbf{Current status:} prototype implementation.
\end{requirementbox}

\begin{futureworkbox}
\textbf{Planned improvements:}
\begin{itemize}
    \item Constant-current source.
    \item RC filter.
    \item Patient protection resistor.
    \item Shielding.
    \item Calibration procedure.
\end{itemize}
\end{futureworkbox}

\subsection{Grounding Strategy}
\label{sec:grounding-strategy}

\begin{engnote}
Each module shall use a single common ground. Ground loops should be avoided, analog return paths kept short, and the microphone return path kept away from high-current power traces where possible.
\end{engnote}

\subsection{Decoupling}
\label{sec:decoupling}

\begin{designnote}
A 0.1\,\textmu F ceramic capacitor is recommended close to each sensor supply, together with a 10\,\textmu F bulk capacitor near the ESP32 power input. Additional bulk capacitance should be added if Wi-Fi brownouts are observed. Some breakout boards already include these components.
\end{designnote}

\subsection{Wiring Recommendations}
\label{sec:wiring-recommendations}

\begin{itemize}
    \item Use 26--28\,AWG stranded wire for power.
    \item Use 28--30\,AWG stranded wire for signal wiring.
    \item Keep I\textsuperscript{2}C runs short.
    \item Use twisted-pair or closely routed GSR leads if cable length increases.
    \item Provide secure strain relief for battery and sensor connections.
\end{itemize}

\subsection{Electrical Safety}
\label{sec:electrical-safety}

\begin{warningbox}
\begin{itemize}
    \item Never apply 5\,V directly to ESP32 GPIO pins.
    \item Verify battery polarity before connecting.
    \item Do not charge the module while connected to GSR electrodes.
    \item Insulate all exposed battery terminals.
    \item Disconnect power before rewiring.
    \item Inspect LiPo cells before every test.
\end{itemize}
\end{warningbox}

\subsection{Estimated Current Budget}
\label{sec:current-budget}

\Cref{tab:current-budget} lists the estimated per-device current draw for both modules. The Wrist Module total is estimated at $\approx$90--210+\,mA, and the Audio Module total at $\approx$85--205+\,mA.

\begin{table}[H]
\centering
\caption{Estimated current budget by module and device.}
\label{tab:current-budget}
\small
\begin{tabularx}{\columnwidth}{@{}l l Y@{}}
\toprule
\textbf{Module} & \textbf{Device} & \textbf{Approx. Current} \\
\midrule
Wrist & ESP32 & 80--200\,mA (Wi-Fi/CPU dependent) \\
Wrist & MAX30102 & 1--5\,mA \\
Wrist & MPU6050 & 3--5\,mA \\
Wrist & ADS1115 & \textless 1\,mA \\
Wrist & MCP9808 & \textless 1\,mA \\
Wrist & GSR front end & Depends on final circuit \\
\midrule
Audio & ESP32 & 80--200\,mA \\
Audio & INMP441 & 1--2\,mA \\
Audio & TEMT6000 & \textless 1\,mA \\
\bottomrule
\end{tabularx}
\end{table}

\subsection{Remaining Electrical Validation}
\label{sec:remaining-electrical-validation}

\begin{requirementbox}
\textbf{The following must still be verified on the assembled hardware:}
\begin{itemize}
    \item Exact ESP32 power input path and acceptable VIN range.
    \item 3.3\,V rail stability under Wi-Fi load.
    \item Battery runtime under continuous operation.
    \item Charge current of the TP4056 (PROG resistor).
    \item Brownout threshold.
    \item Peak current during wireless transmission.
    \item Thermal performance during charging.
    \item Final GPIO validation with firmware.
\end{itemize}
\end{requirementbox}

\subsection{Chapter Status}
\label{sec:ch5-status}

\begin{validationbox}
\textbf{Completed for this chapter:}
\begin{itemize}
    \item Complete wiring architecture
    \item GPIO assignments
    \item I\textsuperscript{2}C wiring
    \item I\textsuperscript{2}S wiring
    \item Power distribution
    \item Grounding strategy
    \item Current budget
    \item Electrical safety
    \item Integration notes
\end{itemize}
\end{validationbox}
